\section{Backend Contract and MuJoCo Physics Realization}
\label{sec:backends}

\subsection{Adapter contract}

The base adapter contract contains seven responsibilities: advertise capabilities,
load a pack and scene, step, snapshot, create a physical connection, remove it, and
shutdown. The snapshot is intentionally small: time, link poses and twists, scalar
joint state, optional measured connector frames, and optional constraint-force
magnitudes. Contacts, solver internals, renderer objects, and engine arrays do not leak
into canonical core state.

Optional structural protocols add joint commands and direct module kinematics. Direct
pose/twist/wrench methods are useful for deterministic scene staging and kinematic
fixtures, but physical demonstrations enforce that these are not called after initial
placement.

\subsection{Lazy registration and selection}

Backends register a name, summary, loader, and installation hint. Resolution is an
explicit argument, then \code{MODSIM\_BACKEND}, then \code{mock}. Loading is lazy, so a
base installation can import the entire core without importing \mujoco{}. A backend
implementation can live in a separate distribution and register through the same
mechanism.

\subsection{Mock backend}

The mock backend is kinematic and dependency-free. It has no gravity, contact, or
articulated link motion. It nonetheless implements the semantic behaviors core needs:
module placement, world-frame twist integration, physical-connection success/failure,
rigid propagation of connected groups, release, injected loads, and snapshot
generation. It is therefore a reference semantics backend, not a low-fidelity physics
claim.

\subsection{MuJoCo scene composition}

The \mujoco{} adapter composes all module placements into one model. For each instance
it:

\begin{enumerate}
  \item selects a same-ID MJCF asset when available, otherwise imports URDF;
  \item attaches the module under a namespace prefix \code{<module>/};
  \item adds a free joint to the module root;
  \item creates connector sites from authored local poses;
  \item creates bounded unit-gear motors for declared effort-controlled scalar joints;
  \item reserves inactive weld and contact-exclusion slots; and
  \item compiles and indexes bodies, joints, sites, actuators, and assets.
\end{enumerate}

The default gravity vector is $(0,0,-9.81)\,\mathrm{m/s^2}$. An optional engine plane is
physically infinite; its visual half-extent is ten meters. Visual, environment, and
debug collision geoms use distinct groups. When importing URDF, the adapter retains
visual geometry unless the asset explicitly requests otherwise. Collision geoms are
placed in hidden group 3 when separate visuals exist, so physics remains active without
opaque proxies obscuring the robot.

\subsection{Time stepping and snapshots}

The backend covers a requested duration with rounded whole solver steps and reports the
actual engine time. Snapshot extraction resolves world body motion, connector site
frames, joint position/velocity, and actuator effort. For anisotropic \smoresep{} tire
contacts, stepping is split so the contact tangent basis can be oriented with the
current wheel axle before constraint solution.

\subsection{Runtime weld pool}

\mujoco{} model topology is compiled, so a runtime dock claims a preallocated inactive
weld rather than adding a new equality object. Each slot has a paired precompiled body
contact exclusion. Claiming a connection sets the equality body operands, relative
transform data, and active flag; it publishes a sorted exclusion signature so the
welded body pair does not simultaneously fight a contact constraint. Release
deactivates and returns both resources.

The connector-relative request must be converted to a body-relative weld:

\begin{equation}
\mathbf T_{B_aB_b}=
\mathbf T_{B_aC_a}
\mathbf T_{C_aC_b}
\mathbf T_{C_bB_b}.
\label{eq:weld-transform}
\end{equation}

The implementation isolates this transformation in a testable module. Nominal snapping
also accounts for connectors on articulated child bodies: it measures root-to-child
transform, solves the desired root pose, clears root and articulated velocity, and only
then activates the weld.

Pool exhaustion and unsupported connection types are ordinary backend refusals. They
propagate to \code{DockFailed} rather than a logical-only success. The current adapter
supports fixed connections only. Compliance, hinge, ball, and custom constraints remain
future work.

\subsection{Effort actuation and state}

The adapter advertises the effort control mode. Scene composition creates one actuator
for each semantically declared effort-controlled joint with a finite effort limit.
Atomic commands resolve all IDs and bounds before writing \code{data.ctrl}. Values
persist until replaced or cleared. Snapshots report position, velocity, and actuator
effort using stable semantic joint IDs.

The architecture can represent position and velocity modes, but the adapter does not
yet implement them. Likewise, the current system does not infer arbitrary transmissions
from URDF. A future actuator catalog should make motor, gear, and controller dynamics
explicit.

\subsection{Native viewer and pacing}

The native passive viewer starts with collision debug group 3 hidden and preserves the
final state for inspection. On macOS, \mujoco{} requires \code{mjpython} for passive
viewer main-thread ownership \citep{mujocoPythonDocs}; the Runtime Inspector resolves a
sibling environment executable automatically.

The launch-time real-time factor $\rho>0$ changes only a wall-clock deadline:

\begin{equation}
t_w^{\mathrm{deadline}} = t_{w,0}+\frac{\Delta t_s}{\rho}.
\label{eq:pacing}
\end{equation}

Simulation timestep, number of steps, motor targets, event timestamps, and simulated
duration remain unchanged. Values below one create slow motion; values above one request
accelerated playback. If physics and rendering cannot keep up, execution proceeds as
fast as possible without skipping physics steps.

At the current head commit, passive viewer synchronization is capped at approximately
30 Hz rather than occurring at every solver step. Synchronization is scheduled from
completion so an expensive render does not trigger another render on the immediately
following physics step, and a final synchronization is guaranteed when the run ends
between refreshes. This decoupling is important at a \SI{2}{ms}-style solver step,
where sync-per-step would attempt 500 visual refreshes per simulated second.

\subsection{Backend comparison}

\begin{table}[htbp]
\centering
\caption{Implemented backend capabilities.}
\label{tab:backend-comparison}
\begin{tabularx}{\textwidth}{L{0.30\textwidth}YY}
\toprule
Capability & Mock & MuJoCo \\
\midrule
Dependencies & Base only & Optional \code{mujoco} and NumPy \\
Dynamics/contact/gravity & No & Yes \\
Articulated kinematics & No & Yes \\
Joint feedback & No & Position, velocity, effort \\
Joint command & No & Effort \\
Measured connector frames & Composed & Connector sites \\
Runtime fixed connection & Yes & Preallocated weld \\
Runtime release & Yes & Weld deactivation \\
Paired contact exclusion & Not applicable & Yes \\
Constraint force reporting & Injectable test value & Not yet \\
Native 3-D viewer & No & Yes \\
\bottomrule
\end{tabularx}
\end{table}

\subsection{Conformance}

Cross-backend tests run a common pack and scenario with gravity disabled and compare
semantic modules, assemblies, connector frames, connection identities, mating
orientation, and committed event sequence. This isolates adapter correctness from
expected physical trajectory differences. The mock is not evidence of a second
physical solver, but it is a valuable executable specification of the semantic
contract.

